% 常用色彩空间之间的变换关系；实线为线性变换，虚线为非线性变换。宽约110mm。
\begin{tikzpicture}[font=\figurefont\small,
  space/.style={draw=FigureBlue,fill=FigureBlueFill,rounded corners=1.8pt,
    minimum width=16mm,minimum height=8mm,align=center,inner sep=1.2mm},
  ref/.style={space,draw=FigureYellow,fill=FigureYellowFill},
  dev/.style={space,draw=FigureRed,fill=FigureRedFill},
  lin/.style={-{Stealth[length=1.8mm]},draw=BookInk,line width=0.8pt},
  nl/.style={-{Stealth[length=1.8mm]},draw=BookInk,dashed,line width=0.8pt},
  lab/.style={font=\figurefont\footnotesize,text=BookMuted}]
  \node[space] (cam) at (0,0) {相机 RGB\\（线性）};
  \node[ref] (xyz) at (2.9,0) {CIE XYZ};
  \node[space] (lrgb) at (5.8,0) {线性 sRGB};
  \node[dev] (srgb) at (8.7,0) {sRGB\\（伽马编码）};
  \node[ref] (lab) at (2.9,-1.7) {CIE Lab};
  \node[dev] (ycc) at (8.7,-1.7) {Y$'$CbCr};
  \node[dev] (hsv) at (5.8,-1.7) {HSV};
  \node[dev] (cmy) at (5.8,1.5) {CMY(K)};
  \draw[lin] (cam) -- node[above,lab] {校正} (xyz);
  \draw[lin] (xyz) -- node[above,lab] {$3\times3$} (lrgb);
  \draw[nl] (lrgb) -- node[above,lab] {编码} (srgb);
  \draw[nl] (xyz) -- node[left,lab] {立方根} (lab);
  \draw[lin] (srgb) -- node[right,lab] {仿射} (ycc);
  \draw[nl] (srgb) -- node[above left,lab] {分段} (hsv);
  \draw[lin] (srgb) |- node[pos=0.75,above,lab] {$1-$} (cmy);
\end{tikzpicture}
